\documentclass[12pt, a4paper]{article}
\usepackage[utf8]{inputenc}
\usepackage[spanish]{babel}
\usepackage{geometry}
\usepackage{graphicx}
\usepackage{hyperref}
\usepackage{listings}
\usepackage{xcolor}
\usepackage{fancyhdr}

\geometry{a4paper, margin=1in}

\hypersetup{
    colorlinks=true,
    linkcolor=blue,
    filecolor=magenta,      
    urlcolor=cyan,
}

\definecolor{codegreen}{rgb}{0,0.6,0}
\definecolor{codegray}{rgb}{0.5,0.5,0.5}
\definecolor{codepurple}{rgb}{0.58,0,0.82}
\definecolor{backcolour}{rgb}{0.95,0.95,0.92}

\lstdefinestyle{kotlinstyle}{
    backgroundcolor=\color{backcolour},   
    commentstyle=\color{codegreen},
    keywordstyle=\color{magenta},
    numberstyle=\tiny\color{codegray},
    stringstyle=\color{codepurple},
    basicstyle=\ttfamily\footnotesize,
    breakatwhitespace=false,         
    breaklines=true,                 
    captionpos=b,                    
    keepspaces=true,                 
    numbers=left,                    
    numbersep=5pt,                  
    showspaces=false,                
    showstringspaces=false,
    showtabs=false,                  
    tabsize=2,
    language=Kotlin
}
\lstset{style=kotlinstyle}

\pagestyle{fancy}
\fancyhf{}
\fancyhead[L]{Informe del Proyecto: MascotaApp}
\fancyfoot[C]{\thepage}

\title{
    \textbf{Informe Técnico del Proyecto}\\
    \vspace{0.5cm}
    \Large MascotaApp
}
\author{Generado por Gemini CLI}
\date{\today}

\begin{document}

\maketitle
\thispagestyle{empty}
\newpage

\tableofcontents
\thispagestyle{empty}
\newpage

\section{Introducción}
El presente documento detalla la estructura, arquitectura y componentes del proyecto de software \textbf{MascotaApp}. Esta aplicación, desarrollada en Kotlin para la plataforma Android, está diseñada para facilitar la gestión, búsqueda y adopción de mascotas.

El objetivo de este informe es proporcionar una visión clara y completa del proyecto, sirviendo como guía para desarrolladores, mantenedores y cualquier persona interesada en comprender la organizaci\"on interna del código fuente. Se analizarán las diferentes capas de la arquitectura, los módulos de funcionalidades y las dependencias clave.

\subsection{Objetivos del Proyecto}
\begin{itemize}
    \item Facilitar la conexión entre dueños de mascotas y personas interesadas en adoptar.
    \item Proveer una plataforma para registrar y mostrar perfiles de mascotas.
    \item Ofrecer funcionalidades de búsqueda, chat y geolocalización.
    \item Garantizar una experiencia de usuario fluida y moderna mediante Jetpack Compose.
\end{itemize}

\newpage

\section{Arquitectura del Software}
El proyecto sigue los principios de la \textbf{Arquitectura Limpia (Clean Architecture)}, separando las responsabilidades en tres capas principales: \texttt{data}, \texttt{domain} y \texttt{ui}. Esta separación promueve un código más mantenible, escalable y testeable.

\begin{figure}[h]
    \centering
    \fbox{
        \begin{BVerbatim}
        /
        |-- data/       (Capa de Datos)
        |-- domain/     (Capa de Dominio)
        `-- ui/         (Capa de Presentación)
        \end{BVerbatim}
    }
    \caption{Estructura de capas principal del proyecto.}
\end{figure}

\subsection{Capa de Dominio (\texttt{domain})}
Esta es la capa central de la aplicación. Contiene la lógica de negocio y es independiente de cualquier framework de UI o base de datos.
\begin{itemize}
    \item \textbf{model}: Define las entidades principales del negocio.
    \begin{itemize}
        \item \texttt{Pet.kt}: Representa el modelo de una mascota.
        \item \texttt{User.kt}: Representa el modelo de un usuario.
    \end{itemize}
    \item \textbf{repository}: Define las interfaces (contratos) que los repositorios de la capa de datos deben implementar. Esto permite que la capa de dominio no dependa directamente de las fuentes de datos.
    \begin{itemize}
        \item \texttt{AuthRepository.kt}
        \item \texttt{PetRepository.kt}
        \item \texttt{UserRepository.kt}
    \end{itemize}
    \item \textbf{usecase}: Contiene los casos de uso, que encapsulan una funcionalidad específica de la aplicación.
    \begin{itemize}
        \item \texttt{Login.kt}
        \item \texttt{Register.kt}
    \end{itemize}
\end{itemize}

\subsection{Capa de Datos (\texttt{data})}
Implementa la lógica de acceso a datos definida en la capa de dominio. Es responsable de obtener datos de fuentes externas como una API REST (Firebase) o una base de datos local.
\begin{itemize}
    \item \textbf{repository}: Implementaciones concretas de las interfaces de repositorio.
    \begin{itemize}
        \item \texttt{AuthRepositoryImpl.kt}
        \item \texttt{PetRepositoryImpl.kt}
        \item \texttt{UserRepositoryImpl.kt}
    \end{itemize}
    \item \textbf{di}: Módulos de inyección de dependencias (probablemente Hilt o Koin).
    \begin{itemize}
        \item \texttt{FirebaseModule.kt}: Provee instancias relacionadas con Firebase (Auth, Firestore).
        \item \texttt{DataModule.kt}: Provee las implementaciones de los repositorios.
    \end{itemize}
\end{itemize}

\subsection{Capa de Presentación (\texttt{ui})}
Contiene todos los componentes relacionados con la interfaz de usuario. Está construida con \textbf{Jetpack Compose}.
\begin{itemize}
    \item \textbf{features}: Cada subdirectorio representa una pantalla o una funcionalidad completa de la UI.
    \item \textbf{navigation}: Gestiona la navegación entre las diferentes pantallas.
    \item \textbf{theme}: Define el estilo visual de la aplicación (colores, tipografía, formas).
\end{itemize}

\newpage

\section{Módulos de la Interfaz de Usuario (\texttt{ui/features})}
La capa de UI está organizada por funcionalidades, donde cada una contiene sus propios componentes, vistas y ViewModels.

\subsection{Autenticación (\texttt{auth})}
Este módulo gestiona el flujo de inicio de sesión, registro y la pantalla de bienvenida.
\begin{itemize}
    \item \textbf{inicio}: Pantalla inicial que da la bienvenida al usuario.
    \item \textbf{login}: Pantalla de inicio de sesión con campos para email/contraseña y opciones de login social.
    \item \textbf{register}: Pantalla de registro de nuevos usuarios.
\end{itemize}

\subsection{Home (\texttt{home})}
Pantalla principal de la aplicación después de que el usuario inicia sesión.
\begin{itemize}
    \item \textbf{view}: \texttt{PrincipalScreen.kt}
    \item \textbf{components}: Contiene una gran cantidad de componentes reutilizables para construir la pantalla, como \texttt{PetCard.kt}, \texttt{CategoriesList.kt}, y \texttt{HeaderSection.kt}.
    \item \textbf{viewmodel}: \texttt{HomeViewModel.kt} para manejar la lógica de la pantalla principal.
\end{itemize}

\subsection{Perfiles de Mascotas (\texttt{perfiles\_mascota})}
Muestra una cuadrícula o lista de perfiles de mascotas disponibles.
\begin{itemize}
    \item \textbf{view}: \texttt{PerfilesMascotaScreen.kt}
    \item \textbf{components}: \texttt{PetGrid.kt}, \texttt{PetCard.kt}.
    \item \textbf{viewmodel}: \texttt{PerfilesMascotaViewModel.kt}.
\end{itemize}

\subsection{Detalle de Mascota (\texttt{perfil\_mascota})}
Muestra la información detallada de una mascota seleccionada.
\begin{itemize}
    \item \textbf{view}: \texttt{PerfilMascotaScreen.kt}
    \item \textbf{components}: \texttt{PetHeader.kt}, \texttt{PetCardInfo.kt}, \texttt{PetOwnerInfo.kt}.
    \item \textbf{viewmodel}: \texttt{PerfilMascotaViewModel.kt}.
\end{itemize}

\subsection{Registro de Mascota (\texttt{Registrar\_mascota})}
Formulario para que los usuarios puedan registrar una nueva mascota en la plataforma.
\begin{itemize}
    \item \textbf{view}: \texttt{RegistrarMascotaScreen.kt}
    \item \textbf{components}: \texttt{PetForm.kt}, \texttt{ImageSelection.kt}.
    \item \textbf{viewmodel}: \texttt{RegistrarMascotaViewModel.kt}.
\end{itemize}

\subsection{Otras Funcionalidades}
La aplicación también incluye otros módulos importantes:
\begin{itemize}
    \item \textbf{buscar}: Funcionalidad de búsqueda.
    \item \textbf{chat}: Sistema de mensajería.
    \item \textbf{Maps}: Integración con mapas para mostrar ubicaciones.
    \item \textbf{notifications}: Centro de notificaciones.
    \item \textbf{perfil\_usuario}: Pantalla de perfil del usuario.
    \item \textbf{configuracion}: Panel de ajustes de la aplicación.
\end{itemize}

\newpage

\section{Navegación}
La navegación de la aplicación se gestiona de forma centralizada en el paquete \texttt{navigation}, utilizando el componente \texttt{NavHost} de Jetpack Compose.

\subsection{\texttt{AppScreens.kt}}
Este archivo define todas las rutas de navegación como una clase sellada (\texttt{sealed class}), lo que proporciona seguridad de tipos y evita errores de tipeo en las rutas.

\begin{lstlisting}[caption={Ejemplo de AppScreens.kt}, label={lst:screens}]
sealed class AppScreens(val route: String) {
    object SplashScreen : AppScreens("splash_screen")
    object InicioScreen : AppScreens("inicio_screen")
    object LoginScreen : AppScreens("login_screen")
    object RegisterScreen : AppScreens("register_screen")
    object PrincipalScreen : AppScreens("principal_screen")
    // ... más pantallas
}
\end{lstlisting}

\subsection{\texttt{AppNavigation.kt}}
Aquí se configura el \texttt{NavHostController} y se define el grafo de navegación, asociando cada ruta de \texttt{AppScreens} con su Composable correspondiente.

\begin{lstlisting}[caption={Ejemplo de AppNavigation.kt}, label={lst:navigation}]
@Composable
fun AppNavigation() {
    val navController = rememberNavController()
    NavHost(
        navController = navController,
        startDestination = AppScreens.SplashScreen.route
    ) {
        composable(AppScreens.SplashScreen.route) {
            CargaScreen(navController)
        }
        composable(AppScreens.InicioScreen.route) {
            InicioScreen(navController)
        }
        composable(AppScreens.LoginScreen.route) {
            LoginScreen(navController)
        }
        // ... más composables
    }
}
\end{lstlisting}

\newpage

\section{Gestión de Dependencias}
El proyecto utiliza inyección de dependencias para desacoplar las clases y facilitar las pruebas. El directorio \texttt{data/di} sugiere el uso de \textbf{Hilt} o \textbf{Koin}.

\subsection{\texttt{FirebaseModule.kt}}
Este módulo es responsable de proveer todas las dependencias relacionadas con Firebase.

\begin{lstlisting}[caption={Ejemplo de FirebaseModule.kt con Hilt}, label={lst:firebase_module}]
@Module
@InstallIn(SingletonComponent::class)
object FirebaseModule {

    @Provides
    @Singleton
    fun provideFirebaseAuth(): FirebaseAuth {
        return FirebaseAuth.getInstance()
    }

    @Provides
    @Singleton
    fun provideFirebaseFirestore(): FirebaseFirestore {
        return FirebaseFirestore.getInstance()
    }
}
\end{lstlisting}

\subsection{\texttt{DataModule.kt}}
Este módulo se encarga de "atar" las interfaces de los repositorios con sus implementaciones concretas.

\begin{lstlisting}[caption={Ejemplo de DataModule.kt con Hilt}, label={lst:data_module}]
@Module
@InstallIn(SingletonComponent::class)
abstract class DataModule {

    @Binds
    @Singleton
    abstract fun bindAuthRepository(
        authRepositoryImpl: AuthRepositoryImpl
    ): AuthRepository

    @Binds
    @Singleton
    abstract fun bindPetRepository(
        petRepositoryImpl: PetRepositoryImpl
    ): PetRepository
        
    // ... otros bindings
}
\end{lstlisting}

\newpage

\section{Estilo y Tema Visual (\texttt{theme})}
La consistencia visual de la aplicación se gestiona en el paquete \texttt{theme}. Este enfoque, estandarizado en proyectos de Jetpack Compose, permite definir y reutilizar fácilmente los elementos de diseño.

\begin{itemize}
    \item \textbf{\texttt{Color.kt}}: Define la paleta de colores de la aplicación, incluyendo colores primarios, secundarios, de fondo, etc. Esto permite cambiar el esquema de colores de toda la app modificando un solo archivo.
    
    \item \textbf{\texttt{Type.kt}}: Contiene la definición de los estilos de tipografía (Typography). Aquí se especifican las fuentes, tamaños y pesos para los diferentes roles de texto como títulos, subtítulos, cuerpo de texto, etc.
    
    \item \textbf{\texttt{Shape.kt}}: Define las formas (Shapes) que se usarán para los componentes, como la redondez de las esquinas en botones, tarjetas y otros elementos de la UI.
    
    \item \textbf{\texttt{Theme.kt}}: Es el archivo central que une todos los elementos anteriores (colores, tipografía y formas) en un único tema de Material Design. El Composable \texttt{MascotaAppTheme} envuelve la aplicación para aplicar este estilo de manera consistente en todos los componentes.
\end{itemize}

\begin{lstlisting}[caption={Uso del Tema en la Aplicación}, label={lst:theme}]
// En MainActivity.kt
class MainActivity : ComponentActivity() {
    override fun onCreate(savedInstanceState: Bundle?) {
        super.onCreate(savedInstanceState)
        setContent {
            MascotaAppTheme { // Aplicando el tema
                Surface(
                    modifier = Modifier.fillMaxSize(),
                    color = MaterialTheme.colorScheme.background
                ) {
                    AppNavigation()
                }
            }
        }
    }
}
\end{lstlisting}

\newpage

\section{Conclusión}
El proyecto \textbf{MascotaApp} está estructurado siguiendo las mejores prácticas de desarrollo de software para Android. La adopción de una Arquitectura Limpia, el uso de Jetpack Compose para la UI y la gestión centralizada de la navegación y dependencias, dan como resultado una base de código robusta, escalable y fácil de mantener.

La clara separación de responsabilidades entre las capas \texttt{data}, \texttt{domain} y \texttt{ui} permite que el desarrollo de nuevas funcionalidades y la corrección de errores se puedan realizar de manera aislada y eficiente. La modularización de la UI por funcionalidades (features) contribuye aún más a esta organización.

Este informe ha proporcionado un mapa detallado de la estructura del proyecto, que servirá como un recurso valioso para el equipo de desarrollo a lo largo del ciclo de vida de la aplicación.

\end{document}
